\subsection{A.1. Conventions and the universal multiplicative answer}

The universal tail comparison, its endpoint, and its binomial products
already occur in Cartan's published account, Section~7,
pp.~269--271 of \cite{A_Cartan1976}. We give the comparison explicitly
to fix the distinction between literal and globally bounded weights,
and then establish the arbitrary-triple and completed formulas.

Throughout this subsection $k$ is an arbitrary commutative unital ring,
$S=\Gamma_k(T)$, and
\[
 I_r=\bigoplus_{w\geq r}kT^{[w]},\qquad
 T^{[a]}T^{[b]}=\binom{a+b}{a}T^{[a+b]}.
\]
Let $\Phi_n=k^{n+1}$ be the simplicial module of vertex functions, with
restriction along maps in $\Delta$, let $T=(1,\ldots,1)$, and let
$\Psi=\Phi/kT$. Put
\[
 A_w^q=\Gamma_k^{w-q}(\Phi)\otimes_k\Lambda_k^q\Psi
 \quad(0\leq q\leq w),\qquad A^q=\bigoplus_{w\geq q}A_w^q.
\]
The differential sends $X^{[a]}$ to $X^{[a-1]}dX$ and has total weight
zero. For a simplicial set $X$ we distinguish three conventions:
\begin{align*}
 C_w^q(X)&=\Hom_{\mathrm{sSet}}(X,A_w^q),&
 C_{\mathrm{gr}}(X)&=\bigoplus_w C_w(X),\\
 C_{\mathrm{lit}}^q(X)&=\Hom_{\mathrm{sSet}}(X,A^q),&
 \widehat C^q(X)&=\Hom_{\mathrm{sSet}}\bigl(X,\prod_w A_w^q\bigr).
\end{align*}
The literal convention permits weights that are finite on each simplex
but unbounded over $X$. It is the convention of scans 7--9 when no
global weight bound is imposed. We write $\Gamma_X$ for the left exact
functor $V\mapsto\Hom_{\mathrm{sSet}}(X,V)$ on simplicial modules.

\paragraph{A-001. Proved: universal resolution and acyclic range.}
For $r\geq0$, $F^rA=\bigoplus_{w\geq r}A_w$ is a resolution of the
constant simplicial module $I_r$. For $q<r$, $F^rA^q$ admits an additive
simplicial contraction and is $\Gamma_X$-acyclic for every $X$.

\emph{Proof.}
In simplicial degree $n$ choose the basis $(T,X_1,\ldots,X_n)$.
Then
\[
 A_n=S\otimes_k\Gamma_k(X_1,\ldots,X_n)
                 \otimes_k\Lambda_k(dX_1,\ldots,dX_n).
\]
The one-variable augmented complex contracts onto its constants by
$h(X^{[a]}dX)=X^{[a+1]}$ and $h(X^{[a]})=0$. Tensoring these
contractions gives a contraction in $n$ variables preserving total
weight, so the constants in $F^r$ are precisely $I_r$. The inclusion
of constants is simplicial, even though the contraction uses a vertex.

The normalized complex of $\Phi$ is the identity disk $k\to k$ in
homological degrees $1,0$. Its additive simplicial contraction remains
a simplicial homotopy after applying $\Gamma^a$ for $a>0$: apply the
functor to each component of the homotopy identities, and use
$\Gamma^a(0)=0$. Tensor with $\Lambda^q\Psi$ and take direct sums.
Every summand of $F^rA^q$ has $a=w-q>0$ when $q<r$.

We spell out the acyclicity consequence. An acyclic simplicial module
$V$ embeds in an injective simplicial module $J$. Every map
$\partial\Delta[n]\to J$ extends across $\Delta[n]$, by injectivity
against the monomorphism of free modules on these simplicial sets.
Thus $J$ is acyclic. The quotient $J/V$ is acyclic, and $J\to J/V$
is a trivial Kan fibration. Maps out of every simplicial set lift
against it; hence $R^1\Gamma_X(V)=0$. Iterating on $J/V$ proves all
higher vanishing. This argument uses acyclicity, and does not identify
it with an additive contraction for arbitrary modules.\hfill$\square$

\paragraph{A-002. Proved: the literal answer as an algebra.}
For every $X,k$ there is a natural isomorphism of cohomologically graded
$S$-algebras
\[
 H^*(C_{\mathrm{lit}}(X))\ \cong\
 \bigoplus_{i\geq0}H^i(X;I_i).
 \tag{A-1}\label{eq:literal}
\]
The product uses ordinary cup product and the coefficient pairing
$I_i\otimes_k I_j\to I_{i+j}$ induced by multiplication in $S$.
The comparison to $H^*(X;S)$ is injective; its degree-$i$ image and
cokernel are respectively
\[
 \ker\bigl(H^i(X;S)\to H^i(X;S/I_i)\bigr),
 \qquad \bigoplus_{w<i}H^i(X;k).
\]

\emph{Proof.}
Dimension shifting in the resolution $F^iA$ gives
\[
 \operatorname{coker}\bigl(\Gamma_XF^iA^{i-1}\to
                   \Gamma_X Z^i(F^iA)\bigr)=R^i\Gamma_X(I_i),
\]
because its terms in degrees $0,\ldots,i-1$ are acyclic by A-001.
For $i=0$ the equality is simply the kernel defining closed forms.
The finite splitting $A=F^iA\oplus\bigoplus_{w<i}A_w$ is preserved
by $\Gamma_X$, and the lower weights have no terms in degrees
$i$ or higher. This proves the additive equality, including the
endpoint $i=w$.

For products, the map of simplicial complexes
$F^iA\otimes_k F^jA\to F^{i+j}A$ lies over the coefficient pairing
$I_i\otimes_k I_j\to I_{i+j}$. The terms and the resolved coefficient
modules are degreewise $k$-flat. The first-quadrant tensor double
complex therefore resolves $I_i\otimes_k I_j$: first take cohomology
in one direction, then in the other. Apply the ordinary-to-derived
sections comparison to this pairing. Its construction using a
cosimplicial acyclic resolution and Alexander--Whitney product makes
the comparison commute with products; equivalently use the
naive-versus-derived cup-product lemma, Stacks Project, Lemma 21.33.1,
tag 0FPJ, opened in this run. Its degree-$i$ edge map is exactly the
dimension-shifting isomorphism above. Thus the displayed product is
the coefficient cup product, with no factorial cancellation.
Finally $S=I_i\oplus\bigoplus_{w<i}kT^{[w]}$ as a $k$-module.
Taking coefficient cohomology proves the image and cokernel formulas.
\hfill$\square$

\paragraph{A-003. Proved: fixed weights and the corrected PD--Rees formula.}
Using the coefficient generator $T^{[w]}$ to normalize the comparison,
\[
 H^i(C_w(X))=\begin{cases}H^i(X;k),&0\leq i\leq w,\\0,&i>w.
 \end{cases}
\]
For $x\in H^i(X;k)$ and $y\in H^j(X;k)$ in allowed weights,
\begin{align*}
 [x]_w[y]_v&=\binom{w+v}{w}[x\smile y]_{w+v},\\
 T^{[a]}[x]_w&=\binom{w+a}{a}[x]_{w+a}.
 \tag{A-2}\label{eq:weights}
\end{align*}
Consequently L1 is true for the genuinely graded cochains:
\[
 H^*(C_{\mathrm{gr}}(X))=
 \bigoplus_{i\geq0}H^i(X;k)\otimes_k I_i.
 \tag{A-3}
\]
\emph{Proof.}
Repeat A-001 and A-002 in one weight. The multiplication covers
$T^{[w]}T^{[v]}=\binom{w+v}{w}T^{[w+v]}$ on constants, proving
both formulas. Direct sums of module complexes are exact, giving
(A-3). The natural map from (A-3) to (A-1) is an isomorphism whenever
cohomology commutes with the relevant countable coefficient sums.
This holds when $X$ has finitely many nondegenerate simplices in each
degree. It also holds when $H_i(X;\mathbb Z)$ and
$H_{i-1}(X;\mathbb Z)$ are finitely generated in each required degree,
by the universal coefficient exact sequence over $\mathbb Z$.
\hfill$\square$

\paragraph{A-004. Refuted: L1 with literal unbounded weights on all $X$.}
Let $k=\mathbb Z$ and let $X=\mathbb N$ be discrete. Then
$H^0(C_{\mathrm{lit}})=\prod_{n\geq0}S$, whereas the proposed tensor
formula maps to it as
\[
 (\prod_n\mathbb Z)\otimes S=\bigoplus_w\prod_n\mathbb Z.
\]
Every element of its image has a uniform finite weight bound. The
sequence $(T^{[n]})_{n\geq0}$ does not. There are no positive-degree
forms on this $X$, so this is the full cohomology calculation and a
complete counterexample. The corrected universal statement is A-002;
the corrected genuinely graded statement is A-003. This infinite
counterexample is proved directly, not by a finite test.

\paragraph{A-005. Proved: universal total-weight completion.}
For all $X,k$,
\[
 H^i(\widehat C(X))=\prod_{w\geq i}H^i(X;k).
 \tag{A-4}
\]
Its product has weight-$n$ component
$\sum_{w+v=n}\binom nw x_w\smile y_v$. Indeed $\Hom$ commutes
with products, products of modules are exact, and each product
coefficient is a finite sum. Thus A-003 proves all claims.

\subsection{A.1 continued. All PD triples and completed coefficients}

Let $B=\Gamma_{\mathbb Z}(T)$ and write $I_r^B$ for its tails. A PD
triple $(S,J,t)$ specifies $f:B\to S$ by
$f(T^{[a]})=\gamma_a(t)$. Its underlying de Rham algebra depends on
this map. Define constant coefficient complexes
\[
 E_r(S)=S\otimes_B^{\mathbf L}I_r^B\in D(S),\qquad E_0=S.
\]
We regard normalized chains $C_*(X;S)$ and $E_r$ as cohomological
complexes in nonpositive degrees and set
\[
 \mathbb H^i(X;E)=\Hom_{D(S)}(C_*(X;S),E[i]).
\]
For a module in degree zero this is ordinary coefficient cohomology.

An entirely explicit free $S$-complex representing $E_r$ is the
normalized two-sided bar complex with degree-$h$ basis
$s[a_1|\cdots|a_h;b]$, where $a_j\geq1$, $b\geq r$. Its
differential for $h>0$ is
\begin{align*}
 \partial(s[a_1|\cdots|a_h;b])={}&
 s\gamma_{a_1}(t)[a_2|\cdots|a_h;b]\\
 &+\sum_{j=1}^{h-1}(-1)^j\binom{a_j+a_{j+1}}{a_j}
 s[a_1|\cdots|a_j+a_{j+1}|\cdots|a_h;b]\\
 &+(-1)^h\binom{a_h+b}{a_h}
 s[a_1|\cdots|a_{h-1};a_h+b].
 \tag{A-5}\label{eq:bar}
\end{align*}
Before tensoring with $S$ this is a $B$-free resolution, since
$B^+$ and $I_r^B$ are free over $\mathbb Z$; the usual extra-degeneracy
bar contraction proves exactness. Thus it works even when $S$ is
not $\mathbb Z$-flat. The product-total Hom complex from normalized
chains into (A-5) computes $\mathbb H$, for arbitrary infinite $X$.

\paragraph{A-006. Proved: the complete answer for arbitrary triples.}
For every PD triple and every simplicial set, $H^0_{\mathrm{DR}}(X)=
H^0(X;S)$, and for $i>0$,
\[
 H^i_{\mathrm{DR}}(X;S,J,t)\cong
 \im\bigl(\mathbb H^i(X;E_i(S))\longrightarrow
              \mathbb H^i(X;E_{i-1}(S))\bigr).
 \tag{A-6}\label{eq:alltriples}
\]
This determines the algebra as follows. Lift image classes in degrees
$i,j$ to $\mathbb H^i(X;E_i)$ and $\mathbb H^j(X;E_j)$, apply
the derived cup product and
$E_i\otimes_S^{\mathbf L}E_j\to E_{i+j}$, and then pass to
the image in $\mathbb H^{i+j}(X;E_{i+j-1})$. This is independent of
the lifts. The coefficient pairing is signed bar shuffle followed by
$I_i^BI_j^B\subset I_{i+j}^B$. Degree zero and $S$ act by coefficient
cup product. These rules include every product, not just the modules.

\emph{Proof.}
Write $Z_B^i=\ker(d:A_B^i\to A_B^{i+1})$. The objectwise
Poincare contraction of A-001 is $B$-linear and splits every exact
differential sequence. Thus each $A_{B,n}^q$ and $Z_{B,n}^q$ is
$B$-projective. Moore normalization is exact and a direct summand
degreewise. Consequently
\[
 NZ_S^i=(NZ_B^i)\otimes_B S,\qquad
 NA_S^q=(NA_B^q)\otimes_B S.
\]
Forms of degree $i$ vanish on simplices of dimension less than $i$.
Dimension shifting in each weight gives the canonical identifications
\[
 NZ_B^i\simeq I_i^B[i],\qquad NA_B^q\simeq\mathbb Z[q]
 \quad\hbox{in }D(B).
 \tag{A-7}
\]
Here the shift places the module in homological degree $i$ or $q$,
and $B$ acts on $\mathbb Z$ by augmentation. For the second equality
all weights $w>q$ are contractible; the sole remaining term
$\Lambda^q\Psi$ has normalization $\mathbb Z[q]$. The binomial
action in A-003 identifies the first module in (A-7) exactly as
$I_i^B$, including its $B$-action. These normalized complexes are
bounded below complexes of projectives and therefore compute derived
tensor products. Hence $NZ_S^i\simeq E_i(S)[i]$.

The degreewise split exact sequence
$0\to Z_S^{i-1}\to A_S^{i-1}\to Z_S^i\to0$ gives a triangle.
Over $B$ it is the shifted rotation of
$0\to I_i^B\to I_{i-1}^B\to\mathbb Z T^{[i-1]}\to0$.
This identification can be checked weightwise: in weights at least
$i$ the connecting map identifies the same constant generator,
while in weight $i-1$ the quotient is the identity. Base change
identifies its last arrow with $E_i[i]\to E_{i-1}[i]$.

The quotient of actual cocycles by actual boundaries is
\[
 \operatorname{coker}\bigl([X,A_S^{i-1}]\to[X,Z_S^i]\bigr).
 \tag{A-8}
\]
Indeed the degreewise epimorphism $A_S^{i-1}\to Z_S^i$ is a Kan
fibration. Every nullhomotopy of a cocycle lifts starting at zero,
so every nullhomotopic cocycle is an actual boundary. Conversely the
image on homotopy classes is exactly the class of actual boundaries.
For a simplicial module $V$ the group $[X,V]$ is
$\Hom_{D(S)}(C_*(X;S),NV)$: linearize and apply Dold--Kan; the
normalized source is bounded below and free. Apply this Hom functor
to the triangle. Exactness identifies (A-8) with (A-6). Closed
zero-forms are constant $S$-valued sections, proving degree zero.

Cycle multiplication and normalized shuffle give, after (A-7), a map
$I_i^B\otimes_B^{\mathbf L}I_j^B\to I_{i+j}^B$. Its source is
homologically connective and its target is in degree zero, so it is
determined by its map on $H_0$. The product comparison A-003 identifies
that map as ordinary ideal multiplication. Tensoring with $S$ gives
the stated pairing on $E$. A boundary times a cycle is a boundary,
by $d(az)=(da)z$ with the appropriate graded sign. This proves lift
independence and equality with the original product.\hfill$\square$

\paragraph{A-007. Proved under the stated topology: completion without
an associated-graded regularity assumption.}
Suppose $S$ is complete and separated for the multiplicative filtration
$F^rS$ of closed PD-tail ideals generated by $\gamma_n(t)$, $n\geq r$.
Complete forms in total weight, giving variables and exterior generators
weight one. Complete (A-5) by assigning its basis element weight
$a_1+\cdots+a_h+b$, in addition to the coefficient filtration. Denote
the resulting complex by $\widehat E_r$. Its degree-$h$ term is the
product of copies of $S$ indexed by these tuples. The first differential
converges since $\gamma_a(t)\in F^aS$; the other contributions to
each coordinate are finite. Then (A-6), with $E$ replaced by
$\widehat E$, and its continuous shuffle product, is the complete
answer for the completed forms.

\emph{Proof.}
The normalized universal cycle complexes and bar complexes are graded
projective resolutions of $I_r^B$. The comparison theorem gives
homogeneous chain maps and homogeneous chain homotopies in both
directions: construct each successive lift against the surjection onto
cycles, using graded projectivity. After coefficient change these
maps preserve total filtration, and hence extend to completions;
their homotopy identities persist. The Poincare splitting is also
filtration preserving, so its split exact cycle sequences remain
exact on completion. Normalization, a degreewise finite-idempotent
operation, commutes with completion. We obtain completed normalized
cycles $\widehat E_r[r]$. The fibration and product arguments in
A-006 now apply to the completed simplicial modules. For the product,
choose before completion the homogeneous chain homotopy identifying
normalized cycle multiplication with the bar-shuffle pairing; both
maps lift the same ideal multiplication between graded projective
resolutions. That homotopy also extends continuously. Thus the
completed product comparison does not assume that a product of
copies of $S$, or the completed bar complex, is $K$-flat.
This proves an actual formula without replacing derived tails by
ordinary ideals.\hfill$\square$

\paragraph{A-008. Proved: the associated-graded hypothesis of scans
233--239.}
Assume in addition that $k=S/F^1S$ and the canonical graded PD map
$\Gamma_k(T)\to\operatorname{gr}_F S$ is an isomorphism. Then
$\widehat E_r\simeq F^rS[0]$, naturally in tails and products.
In particular, for $i>0$,
\[
 H^i_{\widehat{\mathrm{DR}}}(X;S)\cong
 \im\bigl(H^i(X;F^iS)\to H^i(X;F^{i-1}S)\bigr).
 \tag{A-9}
\]
Products are the coefficient cup products on lifts followed by this
image map. Degree zero is $H^0(X;S)$.

\emph{Proof.}
The associated graded of the completed augmented bar complex is the
universal graded $B$-free resolution tensored with
$k\otimes_{\mathbb Z}B$. It is exact in positive homological
degrees, because $I_r^B$ is $\mathbb Z$-free. Its degree-zero
homology is $\bigoplus_{n\geq r}kT^{[n]}=\operatorname{gr}_F F^rS$.
The augmented complex is separated and complete. In degree $h$ its
least basis weight is $r+h$, so below any fixed filtration bound
only finitely many homological degrees occur. Kill a cycle successively
in its least filtration using graded exactness; the successive
primitives converge. The same correction argument proves injectivity
and surjectivity onto $F^rS$ in degree zero. Thus the augmentation
is a quasi-isomorphism, not merely a graded analogy. Naturality
identifies the products and tail maps, and A-007 proves (A-9).
For the split universal completion the tail inclusions split over $k$,
recovering A-005. In a nonsplit base the consecutive-tail image in
(A-9) must be retained; possible connecting homomorphisms are not
discarded.\hfill$\square$

\paragraph{A-009. Refuted: replacing derived tails by ordinary tails
without the hypothesis of A-008.}
Take $S=k=\mathbb F_2$, $J=0$, $t=0$, and
$X=\Delta[2]/\partial\Delta[2]$. In weight two the relative forms,
in coordinates $x=X_1,y=X_2$, are
\[
 k\{xy\}\xrightarrow{d}k\{x\,dy+y\,dx\}
                  \xrightarrow{0}k\{dx\wedge dy\},
 \qquad d(xy)=x\,dy+y\,dx.
\]
The faces $x=0,y=0$ vanish immediately. On the third, $y=-x$,
the first two restrictions vanish since $x^2=2x^{[2]}=0$ and
$2x\,dx=0$. Conversely the first two faces force all other
coefficients to vanish, and the third leaves exactly these spaces.
Thus $H^2$ in weight two is $\mathbb F_2$, although every positive
ordinary PD tail of $t$ is zero. Completion preserves this weight.
The exhaustive calculation is in \texttt{code/A\_t0\_sphere2.py}
and \texttt{logs/A\_t0\_sphere2.json}; the symbolic calculation
just given verifies every hypothesis and all relative generators.
